\section{K}

% kaufen (to buy)
\begin{verbentry}{
  style = {acc},
  toc = {kaufen (to buy)},
  name = {kaufen},
  english = {to buy},
  type = {regular, + Akkusativ},
  auxiliary = {haben}
}
 
    
     \vspace{5pt}

    \hrule
    
    \vspace{5pt}

  \begin{tabular}{rl}
     \multicolumn{2}{l}{\textbf{PRÄSENS} }\\
    ich & \textbf{kaufe}\\
    du & \textbf{kaufst}\\
    er/sie/es & \textbf{kauft}\\
    wir & \textbf{kaufen}\\
    ihr & \textbf{kauft}\\
    sie/Sie & \textbf{kaufen}\\
  \end{tabular}
  \hfill
     \begin{tabular}{rl}
    \multicolumn{2}{l}{\textbf{PRÄTERITUM} }\\
    ich & \textbf{kaufte}\\
    du & \textbf{kauftest}\\
    er/sie/es & \textbf{kaufte}\\
    wir & \textbf{kauften}\\
    ihr & \textbf{kauftet}\\
    sie/Sie & \textbf{kauften}\\
  \end{tabular}
  \hfill
  \begin{tabular}{rl}
     \multicolumn{2}{l}{\textbf{IMPERATIV} }\\
   & \textbf{kauf(e) (du)!}\\
   & \textbf{kaufen wir!}\\
   & \textbf{kauft (ihr)!}\\
   & \textbf{kaufen Sie!}\\
   & \vspace{10pt}
  \end{tabular}

    \vspace{10pt}

 \begin{tabular}{rl}
     \multicolumn{2}{l}{\textbf{PERFEKT}}\\
   & haben + \textbf{gekauft}\\
  \end{tabular}
\hfill
   \begin{tabular}{rl}
   
    \multicolumn{2}{l}{\textbf{FUTURE I} }\\
   & werden + \textbf{kaufen}\\
 
  \end{tabular}
\hfill
   \begin{tabular}{rl}
      \multicolumn{2}{l}{\textbf{PLUSQUAMPERFEKT}}\\
   & hatten + \textbf{gekauft}\\
  \end{tabular}
  
  
\vspace{10pt}

\begin{tabular}{lll}
\multicolumn{3}{l}{\textbf{MIT PRÄPOSITIONEN}}\\
& \textbf{---} & \\
\end{tabular}
\hfill
\begin{tabular}{lll}
\multicolumn{3}{l}{\textbf{TRENNBARE VERBEN}}\\
& \textbf{---} & \\
\end{tabular}
\vspace{-5pt}
\end{verbentry}

% kennen (to know (person/place))
\begin{verbentry}{
  style = {acc},
  toc = {kennen (to know (person/place))},
  name = {kennen},
  english = {to know (person/place)},
  type = {regular, + Akkusativ},
  auxiliary = {haben}
}
 
    
     \vspace{5pt}

    \hrule
    
    \vspace{5pt}

  \begin{tabular}{rl}
     \multicolumn{2}{l}{\textbf{PRÄSENS} }\\
    ich & \textbf{kenne}\\
    du & \textbf{kennst}\\
    er/sie/es & \textbf{kennt}\\
    wir & \textbf{kennen}\\
    ihr & \textbf{kennt}\\
    sie/Sie & \textbf{kennen}\\
  \end{tabular}
  \hfill
     \begin{tabular}{rl}
    \multicolumn{2}{l}{\textbf{PRÄTERITUM} }\\
    ich & \textbf{kannte}\\
    du & \textbf{kanntest}\\
    er/sie/es & \textbf{kannte}\\
    wir & \textbf{kannten}\\
    ihr & \textbf{kanntet}\\
    sie/Sie & \textbf{kannten}\\
  \end{tabular}
  \hfill
  \begin{tabular}{rl}
     \multicolumn{2}{l}{\textbf{IMPERATIV} }\\
   & \textbf{kenn(e) (du)!}\\
   & \textbf{kennen wir!}\\
   & \textbf{kennt (ihr)!}\\
   & \textbf{kennen Sie!}\\
   & \vspace{10pt}
  \end{tabular}

    \vspace{10pt}

 \begin{tabular}{rl}
     \multicolumn{2}{l}{\textbf{PERFEKT}}\\
  &  haben + \textbf{gekannt}\\
  \end{tabular}
\hfill
   \begin{tabular}{rl}
   
    \multicolumn{2}{l}{\textbf{FUTURE I} }\\
  &  werden + \textbf{kennen}\\
 
  \end{tabular}
\hfill
   \begin{tabular}{rl}
      \multicolumn{2}{l}{\textbf{PLUSQUAMPERFEKT}}\\
   & hatten + \textbf{gekannt}\\
  \end{tabular}
  
  
\vspace{10pt}

\begin{tabular}{lll}
\multicolumn{3}{l}{\textbf{MIT PRÄPOSITIONEN}}\\
& \textbf{---} & \\
\end{tabular}
\hfill
\begin{tabular}{lll}
\multicolumn{3}{l}{\textbf{TRENNBARE VERBEN}}\\
& \textbf{---} & \\
\end{tabular}
\vspace{-5pt}
\end{verbentry}

% klären (to clarify)
\begin{verbentry}{
  style = {acc},
  toc = {klären (to clarify)},
  name = {klären},
  english = {to clarify},
  type = {regular, + Akkusativ},
  auxiliary = {haben}
}
 
    
     \vspace{5pt}

    \hrule
    
    \vspace{5pt}

  \begin{tabular}{rl}
     \multicolumn{2}{l}{\textbf{PRÄSENS} }\\
    ich & \textbf{kläre}\\
    du & \textbf{klärst}\\
    er/sie/es & \textbf{klärt}\\
    wir & \textbf{klären}\\
    ihr & \textbf{klärt}\\
    sie/Sie & \textbf{klären}\\
  \end{tabular}
  \hfill
     \begin{tabular}{rl}
    \multicolumn{2}{l}{\textbf{PRÄTERITUM} }\\
    ich & \textbf{klärte}\\
    du & \textbf{klärtest}\\
    er/sie/es & \textbf{klärte}\\
    wir & \textbf{klärten}\\
    ihr & \textbf{klärtet}\\
    sie/Sie & \textbf{klärten}\\
  \end{tabular}
  \hfill
  \begin{tabular}{rl}
     \multicolumn{2}{l}{\textbf{IMPERATIV} }\\
   & \textbf{klär(e) (du)!}\\
   & \textbf{klären wir!}\\
   & \textbf{klärt (ihr)!}\\
   & \textbf{klären Sie!}\\
   & \vspace{10pt}
  \end{tabular}

    \vspace{10pt}

 \begin{tabular}{rl}
     \multicolumn{2}{l}{\textbf{PERFEKT}}\\
   & haben + \textbf{geklärt}\\
  \end{tabular}
\hfill
   \begin{tabular}{rl}
   
    \multicolumn{2}{l}{\textbf{FUTURE I} }\\
   & werden + \textbf{klären}\\
 
  \end{tabular}
\hfill
   \begin{tabular}{rl}
      \multicolumn{2}{l}{\textbf{PLUSQUAMPERFEKT}}\\
   & hatten + \textbf{geklärt}\\
  \end{tabular}
  
  
\vspace{10pt}

\begin{tabular}{lll}
\multicolumn{3}{l}{\textbf{MIT PRÄPOSITIONEN}}\\
& \textbf{---} & \\
\end{tabular}
\hfill
\begin{tabular}{lll}
\multicolumn{3}{l}{\textbf{TRENNBARE VERBEN}}\\
& \textbf{---} & \\
\end{tabular}
\vspace{-5pt}
\end{verbentry}

% kleiden (to dress / to clothe)
\begin{verbentry}{
  style = {default},
  toc = {kleiden (to dress / to clothe)},
  name = {kleiden},
  english = {to dress / to clothe},
  type = {regular},
  auxiliary = {haben}
}
 
    
     \vspace{5pt}

    \hrule
    
    \vspace{5pt}

  \begin{tabular}{rl}
     \multicolumn{2}{l}{\textbf{PRÄSENS} }\\
    ich & \textbf{kleide}\\
    du & \textbf{kleidest}\\
    er/sie/es & \textbf{kleidet}\\
    wir & \textbf{kleiden}\\
    ihr & \textbf{kleidet}\\
    sie/Sie & \textbf{kleiden}\\
  \end{tabular}
  \hfill
     \begin{tabular}{rl}
    \multicolumn{2}{l}{\textbf{PRÄTERITUM} }\\
    ich & \textbf{kleidete}\\
    du & \textbf{kleidetest}\\
    er/sie/es & \textbf{kleidete}\\
    wir & \textbf{kleideten}\\
    ihr & \textbf{kleidetet}\\
    sie/Sie & \textbf{kleideten}\\
  \end{tabular}
  \hfill
  \begin{tabular}{rl}
     \multicolumn{2}{l}{\textbf{IMPERATIV} }\\
   & \textbf{kleide (du)!}\\
   & \textbf{kleiden wir!}\\
   & \textbf{kleidet (ihr)!}\\
   & \textbf{kleiden Sie!}\\
   & \vspace{10pt}
  \end{tabular}

    \vspace{10pt}

 \begin{tabular}{rl}
     \multicolumn{2}{l}{\textbf{PERFEKT}}\\
   & haben + \textbf{gekleidet}\\
  \end{tabular}
\hfill
   \begin{tabular}{rl}
    \multicolumn{2}{l}{\textbf{FUTURE I} }\\
   & werde + \textbf{kleiden}\\
  \end{tabular}
\hfill
   \begin{tabular}{rl}
      \multicolumn{2}{l}{\textbf{PLUSQUAMPERFEKT}}\\
   & hatten + \textbf{gekleidet}\\
  \end{tabular}

   \vspace{10pt}

   \begin{tabular}{lll}
      \multicolumn{3}{l}{\textbf{MIT PRÄPOSITIONEN}}\\
& \textbf{sich kleiden in} + Akk & (to dress in)\\
& \textbf{kleiden mit} + Dat & (to clothe with)\\
\end{tabular}
  \hfill
   \begin{tabular}{lll}
      \multicolumn{3}{l}{\textbf{TRENNBARE VERBEN}}\\
 & \textbf{an$|$kleiden} & (to dress someone)\\
 & \textbf{ver$|$kleiden} & (to disguise / dress up)\\
\vspace{10pt}
  \end{tabular}

  \vspace{-5pt}
\end{verbentry}

% klettern (to climb)
\begin{verbentry}{
  style = {default},
  toc = {klettern (to climb)},
  name = {klettern},
  english = {to climb},
  type = {regular},
  auxiliary = {sein / haben}
}
 
    
     \vspace{5pt}

    \hrule
    
    \vspace{5pt}

  \begin{tabular}{rl}
     \multicolumn{2}{l}{\textbf{PRÄSENS} }\\
    ich & \textbf{klettere}\\
    du & \textbf{kletterst}\\
    er/sie/es & \textbf{klettert}\\
    wir & \textbf{klettern}\\
    ihr & \textbf{klettert}\\
    sie/Sie & \textbf{klettern}\\
  \end{tabular}
  \hfill
     \begin{tabular}{rl}
    \multicolumn{2}{l}{\textbf{PRÄTERITUM} }\\
    ich & \textbf{kletterte}\\
    du & \textbf{klettertest}\\
    er/sie/es & \textbf{kletterte}\\
    wir & \textbf{kletterten}\\
    ihr & \textbf{klettertet}\\
    sie/Sie & \textbf{kletterten}\\
  \end{tabular}
  \hfill
  \begin{tabular}{rl}
     \multicolumn{2}{l}{\textbf{IMPERATIV} }\\
   & \textbf{klettere (du)!}\\
   & \textbf{klettern wir!}\\
   & \textbf{klettert (ihr)!}\\
   & \textbf{klettern Sie!}\\
   & \vspace{10pt}
  \end{tabular}

    \vspace{10pt}

 \begin{tabular}{rl}
     \multicolumn{2}{l}{\textbf{PERFEKT}}\\
   & sein + \textbf{geklettert}\\
  \end{tabular}
\hfill
   \begin{tabular}{rl}
    \multicolumn{2}{l}{\textbf{FUTURE I} }\\
   & werde + \textbf{klettern}\\
  \end{tabular}
\hfill
   \begin{tabular}{rl}
      \multicolumn{2}{l}{\textbf{PLUSQUAMPERFEKT}}\\
   & waren + \textbf{geklettert}\\
  \end{tabular}

   \vspace{10pt}

   \begin{tabular}{lll}
      \multicolumn{3}{l}{\textbf{MIT PRÄPOSITIONEN}}\\
& \textbf{klettern auf} + Akk & (to climb onto)\\
& \textbf{klettern über} + Akk & (to climb over)\\
& \textbf{klettern in} + Akk & (to climb into)\\
\end{tabular}
  \hfill
   \begin{tabular}{lll}
      \multicolumn{3}{l}{\textbf{TRENNBARE VERBEN}}\\
 & \textbf{hinauf$|$klettern} & (to climb up)\\
 & \textbf{herunter$|$klettern} & (to climb down)\\
\vspace{10pt}
  \end{tabular}

  \vspace{-5pt}
\end{verbentry}

% klicken (to click)
\begin{verbentry}{
  style = {default},
  toc = {klicken (to click)},
  name = {klicken},
  english = {to click},
  type = {regular},
  auxiliary = {haben}
}
 

\vspace{5pt}
\hrule
\vspace{5pt}

\begin{tabular}{rl}
\multicolumn{2}{l}{\textbf{PRÄSENS}}\\
ich & \textbf{klicke}\\
du & \textbf{klickst}\\
er/sie/es & \textbf{klickt}\\
wir & \textbf{klicken}\\
ihr & \textbf{klickt}\\
sie/Sie & \textbf{klicken}\\
\end{tabular}
\hfill
\begin{tabular}{rl}
\multicolumn{2}{l}{\textbf{PRÄTERITUM}}\\
ich & \textbf{klickte}\\
du & \textbf{klicktest}\\
er/sie/es & \textbf{klickte}\\
wir & \textbf{klickten}\\
ihr & \textbf{klicktet}\\
sie/Sie & \textbf{klickten}\\
\end{tabular}
\hfill
\begin{tabular}{rl}
\multicolumn{2}{l}{\textbf{IMPERATIV}}\\
& \textbf{klicke (du)!}\\
& \textbf{klicken wir!}\\
& \textbf{klickt (ihr)!}\\
& \textbf{klicken Sie!}\\
\end{tabular}

\vspace{10pt}

\begin{tabular}{rl}
\multicolumn{2}{l}{\textbf{PERFEKT}}\\
& haben + \textbf{geklickt}\\
\end{tabular}
\hfill
\begin{tabular}{rl}
\multicolumn{2}{l}{\textbf{FUTURE I}}\\
& werden + \textbf{klicken}\\
\end{tabular}
\hfill
\begin{tabular}{rl}
\multicolumn{2}{l}{\textbf{PLUSQUAMPERFEKT}}\\
& hatten + \textbf{geklickt}\\
\end{tabular}

\vspace{10pt}

\begin{tabular}{lll}
\multicolumn{3}{l}{\textbf{MIT PRÄPOSITIONEN}}\\
& \textbf{klicken auf} + Akk & (to click on)\\
& \textbf{klicken in} + Akk & (to click into)\\
\end{tabular}
\hfill
\begin{tabular}{lll}
\multicolumn{3}{l}{\textbf{TRENNBARE VERBEN}}\\
& an$|$klicken & (to click on / activate)\\
\vspace{10pt}
\end{tabular}

\vspace{-5pt}
\end{verbentry}

% knien (to kneel)
\begin{verbentry}{
  style = {default},
  toc = {knien (to kneel)},
  name = {knien},
  english = {to kneel},
  type = {regular},
  auxiliary = {haben / sein}
}
 
    
     \vspace{5pt}

    \hrule
    
    \vspace{5pt}

  \begin{tabular}{rl}
     \multicolumn{2}{l}{\textbf{PRÄSENS} }\\
    ich & \textbf{knie}\\
    du & \textbf{kniest}\\
    er/sie/es & \textbf{kniet}\\
    wir & \textbf{knien}\\
    ihr & \textbf{kniet}\\
    sie/Sie & \textbf{knien}\\
  \end{tabular}
  \hfill
     \begin{tabular}{rl}
    \multicolumn{2}{l}{\textbf{PRÄTERITUM} }\\
    ich & \textbf{kniete}\\
    du & \textbf{knietest}\\
    er/sie/es & \textbf{kniete}\\
    wir & \textbf{knieten}\\
    ihr & \textbf{knietet}\\
    sie/Sie & \textbf{knieten}\\
  \end{tabular}
  \hfill
  \begin{tabular}{rl}
     \multicolumn{2}{l}{\textbf{IMPERATIV} }\\
   & \textbf{knie (du)!}\\
   & \textbf{knien wir!}\\
   & \textbf{kniet (ihr)!}\\
   & \textbf{knien Sie!}\\
   & \vspace{10pt}
  \end{tabular}

    \vspace{10pt}

 \begin{tabular}{rl}
     \multicolumn{2}{l}{\textbf{PERFEKT}}\\
  &  haben / sein + \textbf{gekniet}\\
  \end{tabular}
\hfill
   \begin{tabular}{rl}
   
    \multicolumn{2}{l}{\textbf{FUTURE I} }\\
   & werden + \textbf{knien}\\
 
  \end{tabular}
\hfill
   \begin{tabular}{rl}
      \multicolumn{2}{l}{\textbf{PLUSQUAMPERFEKT}}\\
   & hatten / waren + \textbf{gekniet}\\
  \end{tabular}
  
  \vspace{10pt}



\begin{tabular}{lll}
\multicolumn{3}{l}{\textbf{MIT PRÄPOSITIONEN}}\\
& \textbf{---} & \\
\end{tabular}
\hfill
\begin{tabular}{lll}
\multicolumn{3}{l}{\textbf{TRENNBARE VERBEN}}\\
& \textbf{---} & \\
\end{tabular}

  \vspace{-5pt}
\end{verbentry}

% knüpfen (to tie / establish)
\begin{verbentry}{
  style = {default},
  toc = {knüpfen (to tie / establish)},
  name = {knüpfen},
  english = {to tie / establish},
  type = {regular},
  auxiliary = {haben}
}
 

\vspace{5pt}
\hrule
\vspace{5pt}

\begin{tabular}{rl}
\multicolumn{2}{l}{\textbf{PRÄSENS}}\\
ich & \textbf{knüpfe}\\
du & \textbf{knüpfst}\\
er/sie/es & \textbf{knüpft}\\
wir & \textbf{knüpfen}\\
ihr & \textbf{knüpft}\\
sie/Sie & \textbf{knüpfen}\\
\end{tabular}
\hfill
\begin{tabular}{rl}
\multicolumn{2}{l}{\textbf{PRÄTERITUM}}\\
ich & \textbf{knüpfte}\\
du & \textbf{knüpftest}\\
er/sie/es & \textbf{knüpfte}\\
wir & \textbf{knüpften}\\
ihr & \textbf{knüpftet}\\
sie/Sie & \textbf{knüpften}\\
\end{tabular}
\hfill
\begin{tabular}{rl}
\multicolumn{2}{l}{\textbf{IMPERATIV}}\\
& \textbf{knüpfe (du)!}\\
& \textbf{knüpfen wir!}\\
& \textbf{knüpft (ihr)!}\\
& \textbf{knüpfen Sie!}\\
\end{tabular}

\vspace{10pt}

\begin{tabular}{rl}
\multicolumn{2}{l}{\textbf{PERFEKT}}\\
& haben + \textbf{geknüpft}\\
\end{tabular}
\hfill
\begin{tabular}{rl}
\multicolumn{2}{l}{\textbf{FUTURE I}}\\
& werden + \textbf{knüpfen}\\
\end{tabular}
\hfill
\begin{tabular}{rl}
\multicolumn{2}{l}{\textbf{PLUSQUAMPERFEKT}}\\
& hatten + \textbf{geknüpft}\\
\end{tabular}

\vspace{10pt}

\begin{tabular}{lll}
\multicolumn{3}{l}{\textbf{MIT PRÄPOSITIONEN}}\\
& \textbf{knüpfen an} + Akk & (to link to)\\
\end{tabular}
\hfill
\begin{tabular}{lll}
\multicolumn{3}{l}{\textbf{TRENNBARE VERBEN}}\\
& \textbf{an$|$knüpfen} & (to connect to)\\
\end{tabular}

\vspace{-5pt}
\end{verbentry}

% kochen (to cook)
\begin{verbentry}{
  style = {acc},
  toc = {kochen (to cook)},
  name = {kochen},
  english = {to cook},
  type = {regular, + Akkusativ},
  auxiliary = {haben}
}
 
    
     \vspace{5pt}

    \hrule
    
    \vspace{5pt}

  \begin{tabular}{rl}
     \multicolumn{2}{l}{\textbf{PRÄSENS} }\\
    ich & \textbf{koche}\\
    du & \textbf{kochst}\\
    er/sie/es & \textbf{kocht}\\
    wir & \textbf{kochen}\\
    ihr & \textbf{kocht}\\
    sie/Sie & \textbf{kochen}\\
  \end{tabular}
  \hfill
     \begin{tabular}{rl}
    \multicolumn{2}{l}{\textbf{PRÄTERITUM} }\\
    ich & \textbf{kochte}\\
    du & \textbf{kochtest}\\
    er/sie/es & \textbf{kochte}\\
    wir & \textbf{kochten}\\
    ihr & \textbf{kochtet}\\
    sie/Sie & \textbf{kochten}\\
  \end{tabular}
  \hfill
  \begin{tabular}{rl}
     \multicolumn{2}{l}{\textbf{IMPERATIV} }\\
   & \textbf{koch(e) (du)!}\\
   & \textbf{kochen wir!}\\
   & \textbf{kocht (ihr)!}\\
   & \textbf{kochen Sie!}\\
   & \vspace{10pt}
  \end{tabular}

    \vspace{10pt}

 \begin{tabular}{rl}
     \multicolumn{2}{l}{\textbf{PERFEKT}}\\
   & haben + \textbf{gekocht}\\
  \end{tabular}
\hfill
   \begin{tabular}{rl}
   
    \multicolumn{2}{l}{\textbf{FUTURE I} }\\
   & werden + \textbf{kochen}\\
 
  \end{tabular}
\hfill
   \begin{tabular}{rl}
      \multicolumn{2}{l}{\textbf{PLUSQUAMPERFEKT}}\\
    & hatten + \textbf{gekocht}\\
  \end{tabular}
  
  
\vspace{10pt}

\begin{tabular}{lll}
\multicolumn{3}{l}{\textbf{MIT PRÄPOSITIONEN}}\\
& \textbf{---} & \\
\end{tabular}
\hfill
\begin{tabular}{lll}
\multicolumn{3}{l}{\textbf{TRENNBARE VERBEN}}\\
& \textbf{---} & \\
\end{tabular}
\vspace{-5pt}
\end{verbentry}

% kommen (to come)
\begin{verbentry}{
  style = {default},
  toc = {kommen (to come)},
  name = {kommen},
  english = {to come},
  type = {irregular},
  auxiliary = {sein}
}
 
    
     \vspace{5pt}

    \hrule
    
    \vspace{5pt}

  \begin{tabular}{rl}
     \multicolumn{2}{l}{\textbf{PRÄSENS} }\\
    ich & \textbf{komme}\\
    du & \textbf{kommst}\\
    er/sie/es & \textbf{kommt}\\
    wir & \textbf{kommen}\\
    ihr & \textbf{kommt}\\
    sie/Sie & \textbf{kommen}\\
  \end{tabular}
  \hfill
     \begin{tabular}{rl}
    \multicolumn{2}{l}{\textbf{PRÄTERITUM} }\\
    ich & \textbf{kam}\\
    du & \textbf{kamst}\\
    er/sie/es & \textbf{kam}\\
    wir & \textbf{kamen}\\
    ihr & \textbf{kamt}\\
    sie/Sie & \textbf{kamen}\\
  \end{tabular}
  \hfill
  \begin{tabular}{rl}
     \multicolumn{2}{l}{\textbf{IMPERATIV} }\\
   & \textbf{komm(e) (du)!}\\
   & \textbf{kommen wir!}\\
   & \textbf{kommt (ihr)!}\\
   & \textbf{kommen Sie!}\\
   & \vspace{10pt}
  \end{tabular}

    \vspace{10pt}

 \begin{tabular}{rl}
     \multicolumn{2}{l}{\textbf{PERFEKT}}\\
   & sein + \textbf{gekommen}\\
  \end{tabular}
\hfill
   \begin{tabular}{rl}
   
    \multicolumn{2}{l}{\textbf{FUTURE I} }\\
   & werden + \textbf{kommen}\\
 
  \end{tabular}
\hfill
   \begin{tabular}{rl}
      \multicolumn{2}{l}{\textbf{PLUSQUAMPERFEKT}}\\
   & waren + \textbf{gekommen}\\
  \end{tabular}

  \vspace{10pt}

   \begin{tabular}{lll}
      \multicolumn{3}{l}{\textbf{MIT PRÄPOSITIONEN}}\\
& \textbf{kommen aus} + Dat & (to come from)\\
& \textbf{kommen von} + Dat & (origin / source of)\\
&\textbf{kommen zu} + Dat & (to approach to)\\
&\textbf{kommen nach} + Dat & (come to a city / country)\\
&\textbf{kommen in} + Akk & (to come into)\\
\vspace{50pt}
  \end{tabular}
  \hfill
   \begin{tabular}{lll}
      \multicolumn{3}{l}{\textbf{TRENNBARE VERBEN}}\\
 & \textbf{an$|$kommen} & (to arrive)\\
 & \textbf{mit$|$kommen} & (to come along)\\
     & \textbf{auf$|$kommen} & (to arise / come up)\\
     & \textbf{vor$|$kommen} & (to occur)\\
     & \textbf{nach$|$kommen} & (to follow / comply)\\
     & \textbf{zu$|$kommen} & (to approach)\\
      & \textbf{herein$|$kommen} & (to come in)\\
     & \textbf{hinein$|$kommen} & (to get inside / enter)\\
     & \textbf{dazu$|$kommen} & (to join / be added)\\
      & \textbf{fort$|$kommen} & (to make progress)
  \end{tabular}

  \vspace{-5pt}
\end{verbentry}

% kommunizieren (to communicate)
\begin{verbentry}{
  style = {default},
  toc = {kommunizieren (to communicate)},
  name = {kommunizieren},
  english = {to communicate},
  type = {regular},
  auxiliary = {haben}
}
 

\vspace{5pt}
\hrule
\vspace{5pt}

\begin{tabular}{rl}
\multicolumn{2}{l}{\textbf{PRÄSENS}}\\
ich & \textbf{kommuniziere}\\
du & \textbf{kommunizierst}\\
er/sie/es & \textbf{kommuniziert}\\
wir & \textbf{kommunizieren}\\
ihr & \textbf{kommuniziert}\\
sie/Sie & \textbf{kommunizieren}\\
\end{tabular}
\hfill
\begin{tabular}{rl}
\multicolumn{2}{l}{\textbf{PRÄTERITUM}}\\
ich & \textbf{kommunizierte}\\
du & \textbf{kommuniziertest}\\
er/sie/es & \textbf{kommunizierte}\\
wir & \textbf{kommunizierten}\\
ihr & \textbf{kommuniziertet}\\
sie/Sie & \textbf{kommunizierten}\\
\end{tabular}
\hfill
\begin{tabular}{rl}
\multicolumn{2}{l}{\textbf{IMPERATIV}}\\
& \textbf{kommuniziere (du)!}\\
& \textbf{kommunizieren wir!}\\
& \textbf{kommuniziert (ihr)!}\\
& \textbf{kommunizieren Sie!}\\
\end{tabular}

\vspace{10pt}

\begin{tabular}{rl}
\multicolumn{2}{l}{\textbf{PERFEKT}}\\
& haben + \textbf{kommuniziert}\\
\end{tabular}
\hfill
\begin{tabular}{rl}
\multicolumn{2}{l}{\textbf{FUTURE I}}\\
& werden + \textbf{kommunizieren}\\
\end{tabular}
\hfill
\begin{tabular}{rl}
\multicolumn{2}{l}{\textbf{PLUSQUAMPERFEKT}}\\
& hatten + \textbf{kommuniziert}\\
\end{tabular}

\vspace{10pt}

\begin{tabular}{lll}
\multicolumn{3}{l}{\textbf{MIT PRÄPOSITIONEN}}\\
& \textbf{kommunizieren mit} + Dat & (to communicate with)\\
& \textbf{kommunizieren über} + Akk & (to communicate about)\\
\end{tabular}
\hfill
\begin{tabular}{lll}
\multicolumn{3}{l}{\textbf{TRENNBARE VERBEN}}\\
& \textbf{---} & \\
\end{tabular}

\vspace{-5pt}
\end{verbentry}

% kosten (to cost)
\begin{verbentry}{
  style = {acc},
  toc = {kosten (to cost)},
  name = {kosten},
  english = {to cost},
  type = {regular, + Akkusativ},
  auxiliary = {haben}
}
 
    
     \vspace{5pt}

    \hrule
    
    \vspace{5pt}

  \begin{tabular}{rl}
     \multicolumn{2}{l}{\textbf{PRÄSENS} }\\
    ich & \textbf{-}\\
    du & \textbf{-}\\
    er/sie/es & \textbf{kostet}\\
    wir & \textbf{-}\\
    ihr & \textbf{-}\\
    sie/Sie & \textbf{kosten}\\
  \end{tabular}
  \hfill
     \begin{tabular}{rl}
    \multicolumn{2}{l}{\textbf{PRÄTERITUM} }\\
    ich & \textbf{-}\\
    du & \textbf{-}\\
    er/sie/es & \textbf{kostete}\\
    wir & \textbf{-}\\
    ihr & \textbf{-}\\
    sie/Sie & \textbf{kosteten}\\
  \end{tabular}
  \hfill
  \begin{tabular}{rl}
     \multicolumn{2}{l}{\textbf{IMPERATIV} }\\
   & \textbf{-}\\
   & \textbf{-}\\
   & \textbf{-}\\
   & \textbf{-}\\
   & \vspace{10pt}
  \end{tabular}

    \vspace{10pt}

 \begin{tabular}{rl}
     \multicolumn{2}{l}{\textbf{PERFEKT}}\\
   & haben + \textbf{gekostet}\\
  \end{tabular}
\hfill
   \begin{tabular}{rl}
   
    \multicolumn{2}{l}{\textbf{FUTURE I} }\\
   & werden + \textbf{kosten}\\
 
  \end{tabular}
\hfill
   \begin{tabular}{rl}
      \multicolumn{2}{l}{\textbf{PLUSQUAMPERFEKT}}\\
   & hatten + \textbf{gekostet}\\
  \end{tabular}
  
  
\vspace{10pt}

\begin{tabular}{lll}
\multicolumn{3}{l}{\textbf{MIT PRÄPOSITIONEN}}\\
& \textbf{---} & \\
\end{tabular}
\hfill
\begin{tabular}{lll}
\multicolumn{3}{l}{\textbf{TRENNBARE VERBEN}}\\
& \textbf{---} & \\
\end{tabular}
\vspace{-5pt}
\end{verbentry}

% kucken (to look)
\begin{verbentry}{
  style = {default},
  toc = {kucken (to look)},
  name = {kucken},
  english = {to look},
  type = {regular},
  auxiliary = {haben}
}
 
    
     \vspace{5pt}

    \hrule
    
    \vspace{5pt}

  \begin{tabular}{rl}
     \multicolumn{2}{l}{\textbf{PRÄSENS} }\\
    ich & \textbf{kucke}\\
    du & \textbf{kuckst}\\
    er/sie/es & \textbf{kuckt}\\
    wir & \textbf{kucken}\\
    ihr & \textbf{kuckt}\\
    sie/Sie & \textbf{kucken}\\
  \end{tabular}
  \hfill
     \begin{tabular}{rl}
    \multicolumn{2}{l}{\textbf{PRÄTERITUM} }\\
    ich & \textbf{kuckte}\\
    du & \textbf{kucktest}\\
    er/sie/es & \textbf{kuckte}\\
    wir & \textbf{kuckten}\\
    ihr & \textbf{kucktet}\\
    sie/Sie & \textbf{kuckten}\\
  \end{tabular}
  \hfill
  \begin{tabular}{rl}
     \multicolumn{2}{l}{\textbf{IMPERATIV} }\\
   & \textbf{kuck(e) du!}\\
   & \textbf{kucken wir!}\\
   & \textbf{kuckt (ihr)!}\\
   & \textbf{kucken Sie!}\\
   & \vspace{10pt}
  \end{tabular}

    \vspace{10pt}

 \begin{tabular}{rl}
     \multicolumn{2}{l}{\textbf{PERFEKT}}\\
   & haben + \textbf{gekuckt}\\
  \end{tabular}
\hfill
   \begin{tabular}{rl}
   
    \multicolumn{2}{l}{\textbf{FUTURE I} }\\
   & werden + \textbf{kucken}\\
 
  \end{tabular}
\hfill
   \begin{tabular}{rl}
      \multicolumn{2}{l}{\textbf{PLUSQUAMPERFEKT}}\\
   & hatten + \textbf{gekuckt}\\
  \end{tabular}
  
  \vspace{10pt}



\begin{tabular}{lll}
\multicolumn{3}{l}{\textbf{MIT PRÄPOSITIONEN}}\\
& \textbf{---} & \\
\end{tabular}
\hfill
\begin{tabular}{lll}
\multicolumn{3}{l}{\textbf{TRENNBARE VERBEN}}\\
& \textbf{---} & \\
\end{tabular}

  \vspace{-5pt}
\end{verbentry}

% kündigen (to quit / to terminate)
\begin{verbentry}{
  style = {default},
  toc = {kündigen (to quit / to terminate)},
  name = {kündigen},
  english = {to quit / to terminate},
  type = {regular},
  auxiliary = {haben}
}
 
    
     \vspace{5pt}

    \hrule
    
    \vspace{5pt}

  \begin{tabular}{rl}
     \multicolumn{2}{l}{\textbf{PRÄSENS} }\\
    ich & \textbf{kündige}\\
    du & \textbf{kündigst}\\
    er/sie/es & \textbf{kündigt}\\
    wir & \textbf{kündigen}\\
    ihr & \textbf{kündigt}\\
    sie/Sie & \textbf{kündigen}\\
  \end{tabular}
  \hfill
     \begin{tabular}{rl}
    \multicolumn{2}{l}{\textbf{PRÄTERITUM} }\\
    ich & \textbf{kündigte}\\
    du & \textbf{kündigtest}\\
    er/sie/es & \textbf{kündigte}\\
    wir & \textbf{kündigten}\\
    ihr & \textbf{kündigtet}\\
    sie/Sie & \textbf{kündigten}\\
  \end{tabular}
  \hfill
  \begin{tabular}{rl}
     \multicolumn{2}{l}{\textbf{IMPERATIV} }\\
   & \textbf{kündige (du)!}\\
   & \textbf{kündigen wir!}\\
   & \textbf{kündigt (ihr)!}\\
   & \textbf{kündigen Sie!}\\
   & \vspace{10pt}
  \end{tabular}

    \vspace{10pt}

 \begin{tabular}{rl}
     \multicolumn{2}{l}{\textbf{PERFEKT}}\\
   & haben + \textbf{gekündigt}\\
  \end{tabular}
\hfill
   \begin{tabular}{rl}
    \multicolumn{2}{l}{\textbf{FUTURE I} }\\
   & werde + \textbf{kündigen}\\
  \end{tabular}
\hfill
   \begin{tabular}{rl}
      \multicolumn{2}{l}{\textbf{PLUSQUAMPERFEKT}}\\
   & hatten + \textbf{gekündigt}\\
  \end{tabular}

   \vspace{10pt}

   \begin{tabular}{lll}
      \multicolumn{3}{l}{\textbf{MIT PRÄPOSITIONEN}}\\
& \textbf{kündigen bei} + Dat & (to quit at / resign from)\\
& \textbf{kündigen wegen} + Gen & (to terminate because of)\\
\end{tabular}
  \hfill
   \begin{tabular}{lll}
      \multicolumn{3}{l}{\textbf{TRENNBARE VERBEN}}\\
 & \textbf{an$|$kündigen} & (to announce)\\
 & \textbf{auf$|$kündigen} & (to terminate / cancel)\\
\vspace{10pt}
  \end{tabular}

  \vspace{-5pt}
\end{verbentry}
